\scanpage{425}
\mspage{337}

Il reste seulement encore à prouver l'implication
\[
\deleted{a}\unsure{b'|a'}\Longrightarrow a,
\]
ce qui est essentiellement trivial.

Cela achève la démonstration du critère \deleted{\unsure{de}}d'accessibilité \inserted{4.13.15} d'une sous-catégorie pleine $\underline S$ de $\unsure{M|\mathcal M}$.

\deleted{\uline{Corollaire} 4.13.18. Soient $M$ une catégorie \inserted{accessible}, \deleted{\unsure{un cardinal} tel \unsure{$\pi_0$}\inserted{\unsure{$\pi_0$-accessible}} que $M$ \unsure{satisfasse|satisfait} $\mathrm L_{\pi_0}$ ($M$ stable \unsure{$\pi$} \unsure{(un cardinal} \unsure{est} utile pour $M$, par $\varinjlim$ \unsure{grdes|grds|strictes} \unsure{devant|sur} $\pi_0$)}, $\pi$ un cardinal utile pour $M$, $\underline S$ une sous-catégorie str. pleine de $M$, $\underline S(\pi)=\underline S\cap M_\pi$. Soit $\alpha_0$ un ordinal limite, tel que $\mathrm{card}\,\alpha_0\leq\pi$, \deleted{et tel que $I_{\alpha_0}\unsure{:=|\!=}\unsure{\alpha<\alpha_0|[\alpha<\alpha_0]}$ soit \unsure{grd} devant $\pi_0$.} \deleted{et désignons} et tel \deleted{appelons \unsure{$M$} $\mathrm L(\unsure{\alpha|\vec\alpha})$ l'hypothèse : $\underline S(\pi)$ est} stable dans $M_\pi$ par $\varinjlim$ de type $I_{\alpha_0}$. (NB $M_{\unsure{\pi}}$ est stable par ces limites, et $M_\pi$ est stable dans $M$ pour ces limites). \deleted{\unsure{Donc} que $\underline S$}}

On va donner une variante un peu \unsure{décantée|décanté|différente} de 4.13.15, qui (avec des hyp. \unsure{un peu} plus faibles) capte l'essentiel de sa substance. L'énoncé 4.13.15 en est une conséquence \unsure{presque} immédiate :

% archive sheet number: 424

